\section{Method}
\label{sec:method}

\subsection{Dataset and split}
The experiments use the supplied product interaction data, product metadata, and warm-item candidate list. The training file contains 427{,}957 user--item rating interactions from 54{,}483 users. The candidate set contains 22{,}753 unique warm items and matches the item universe available in training. Product metadata provides item titles and the \texttt{content\_text\_basic} field for content representation. We use the provided training and validation split: model parameters and item popularity are derived from training interactions only, while validation interactions are reserved for ranking evaluation and parameter selection. The train and validation user--item pairs have no overlap in the supplied audit. The evaluated cohort consists of 11{,}850 known training users with at least one relevant warm-item validation target.

\subsection{Candidate generation and seen-item filtering}
All methods rank the same warm candidate set. Before ranking for user $u$, items with which $u$ interacted in training are removed. Each method returns up to $K=10$ unique items. Validation items outside the warm candidate set are not rankable and are excluded from the relevant target set. This common candidate and filtering policy is used for the popularity, matrix-factorization, content-based, and reranked lists.

\subsection{Relevance and long-tail definitions}
A validation interaction is relevant when its explicit rating is at least 4. Item popularity is the number of interactions for the item in the training data. We define the long tail as items with at most five training interactions. In addition to long-tail exposure, we report Long-Tail Recall@10, defined as the number of relevant validation user--item pairs whose item appears in that user's Top-10, restricted to long-tail items, divided by all relevant validation long-tail user--item pairs in the evaluated cohort. There are 4{,}599 such target pairs under this protocol.

\subsection{Recommendation methods}
The popularity baseline ranks candidate items by training interaction count. The collaborative baseline applies rank-32 Truncated SVD to the sparse user--item matrix of observed explicit ratings (1--5); missing entries are represented as zeros in this matrix. The content-based baseline creates TF--IDF item vectors from product title and \texttt{content\_text\_basic}, then scores items by cosine similarity to a user profile built from that user's training interactions. These methods are reported as separate baselines.

For the principal popularity-aware experiment, let $s_{ui}$ be the original MF score and let $\widetilde{s}_{ui}$ be its z-score over candidates available to user $u$. Let $p_i$ be the training interaction count and $p_{\max}$ its maximum over candidates. The reranking score is
\begin{equation}
 r_{ui}=\widetilde{s}_{ui}-\lambda\frac{\log(1+p_i)}{\log(1+p_{\max})}+\beta\,\mathbb{I}[p_i\leq 5].
\end{equation}
The retained main configuration uses $\lambda=1$ and $\beta=6$. Controlled validation sweeps vary the popularity penalty and long-tail bonus; an additional sweep gates the long-tail bonus to items within specified original MF rank windows. An observed-only explicit MF objective is also evaluated experimentally. It minimizes squared error only on observed training ratings with L2 regularization. Its 1,000-iteration fit did not converge and is not used to replace the principal baseline.

\subsection{Evaluation metrics}
We report NDCG@10 for binary relevance at rating $\geq4$, catalog coverage (unique recommended items divided by the candidate catalog size), average training popularity of recommended slots, and long-tail exposure (fraction of recommendation slots occupied by items with training count $\leq5$). We additionally report the number of relevant long-tail targets retrieved and Long-Tail Recall@10. Parameter sweeps use validation and therefore are not independent final estimates. The supplied data contain no separate test set; independent test evaluation remains future work.
